\documentclass[11pt]{article}

% ============================================================
% GUÍA Y RÚBRICA — Práctica fuera de clases (PFCl)
% Modelado de casos de uso en herramienta CASE y documentación
% de user stories con criterios de aceptación
% Ingeniería de Requisitos (ISR-401) · UTEQ · 2026-2027 PPA
% ============================================================

\usepackage[utf8]{inputenc}
\usepackage[T1]{fontenc}
\usepackage[spanish]{babel}
\usepackage[a4paper,top=2cm,bottom=2.2cm,left=1.6cm,right=1.6cm]{geometry}
\usepackage{array}
\usepackage{tabularx}
\usepackage{longtable}
\usepackage{pdflscape}
\usepackage{colortbl}
\usepackage{xcolor}
\usepackage{enumitem}
\usepackage{fancyhdr}
\usepackage{lastpage}
\usepackage{tcolorbox}
\usepackage{booktabs}
\usepackage{amsmath}
\usepackage{titlesec}
\usepackage[hidelinks]{hyperref}

% ---------- Paleta institucional ----------
\definecolor{utegreen}{HTML}{0B7A3B}
\definecolor{utegreendk}{HTML}{075C2C}
\definecolor{midblue}{HTML}{1F5C8B}
\definecolor{gold}{HTML}{D4A017}
\definecolor{rowgray}{HTML}{EFEFEF}
\definecolor{boxgray}{HTML}{F4F6F4}
\definecolor{lightred}{HTML}{FBEAEA}
\definecolor{lightgold}{HTML}{FBF3DC}
\definecolor{lightblue}{HTML}{EAF1F7}
\definecolor{lvlexc}{HTML}{DCEFE3}
\definecolor{lvlsat}{HTML}{EAF1F7}
\definecolor{lvldev}{HTML}{FBF3E2}
\definecolor{lvlins}{HTML}{F7E6E6}

% ---------- Columnas personalizadas ----------
\newcolumntype{L}[1]{>{\raggedright\arraybackslash}p{#1}}
\newcolumntype{C}[1]{>{\centering\arraybackslash}p{#1}}

% ---------- Encabezado / pie ----------
\setlength{\headheight}{30pt}
\pagestyle{fancy}
\fancyhf{}
\lhead{\footnotesize\textbf{Ingeniería de Requisitos (ISR-401)}}
\rhead{\footnotesize UTEQ \,\textcolor{gold}{$\bullet$}\, 2026--2027 PPA}
\lfoot{\footnotesize Guía y Rúbrica --- Práctica fuera de clases: Casos de uso y user stories}
\rfoot{\footnotesize Página \thepage\ de \pageref{LastPage}}
\renewcommand{\headrulewidth}{0.6pt}
\renewcommand{\footrulewidth}{0.4pt}

% ---------- Atajos ----------
\newcommand{\thdr}[1]{\cellcolor{utegreen}\textcolor{white}{\textbf{#1}}}
\newcommand{\bhdr}[1]{\cellcolor{midblue}\textcolor{white}{\textbf{#1}}}

\renewcommand{\arraystretch}{1.25}

\begin{document}

% ============================================================
% PORTADA / TÍTULO
% ============================================================
\begin{center}
{\large\textbf{\textcolor{utegreendk}{UNIVERSIDAD TÉCNICA ESTATAL DE QUEVEDO}}}\\[2pt]
{\small Facultad de Ciencias de la Computación \,\textcolor{gold}{$\bullet$}\, Carrera de Software}\\[8pt]
{\LARGE\textbf{\textcolor{utegreendk}{Guía y Rúbrica de Evaluación}}}\\[4pt]
{\large Práctica fuera de clases: \textit{Modelado de casos de uso en herramienta CASE y documentación de user stories con criterios de aceptación}}\\[2pt]
{\small Actividad vinculada al Proyecto Fin de Curso (PFC)}
\end{center}

\vspace{4pt}
\noindent\textcolor{gold}{\rule{\linewidth}{1.5pt}}

% ============================================================
% 1. DATOS GENERALES
% ============================================================
\section*{\textcolor{utegreendk}{1. DATOS GENERALES DE LA ACTIVIDAD}}

\noindent\begin{tabularx}{\linewidth}{|L{0.28\linewidth}|X|}
\hline
\multicolumn{2}{|l|}{\cellcolor{utegreen}\textcolor{white}{\textbf{DATOS GENERALES DE LA ACTIVIDAD}}}\\
\hline
\textbf{Asignatura} & Ingeniería de Requisitos (ISR-401) -- 4.º nivel, Carrera de Software\\
\hline
\rowcolor{rowgray}\textbf{Unidad / temas} & Unidad 2 (Tema 7: metas, escenarios y casos de uso) con proyección a la Unidad 3 (modelos de especificación UML) y al Tema 14 (user stories y criterios de aceptación)\\
\hline
\textbf{Resultado de aprendizaje} & Analiza y modela conceptualmente los requisitos mediante diagramas UML (casos de uso) y escenarios ágiles (user stories), asegurando coherencia, trazabilidad y verificabilidad.\\
\hline
\rowcolor{rowgray}\textbf{Actividad evaluada} & Construcción del \textbf{modelo de casos de uso del sistema del PFC} en una herramienta CASE (\textbf{Enterprise Architect} o \textbf{StarUML}) y documentación de \textbf{tres (3) user stories} con sus \textbf{criterios de aceptación} verificables.\\
\hline
\textbf{Modalidad} & Trabajo en equipo (equipos del PFC). Práctica fuera de clases (trabajo autónomo).\\
\hline
\rowcolor{rowgray}\textbf{Producto entregable} & Informe en PDF subido al repositorio GitHub del equipo, acompañado del \textbf{archivo fuente nativo} del proyecto CASE (\texttt{.qea}/\texttt{.eap}/\texttt{.eapx} para Enterprise Architect; \texttt{.mdj} para StarUML).\\
\hline
\textbf{Referentes} & ISO/IEC/IEEE 29148:2018 (caracterís\-ticas de calidad de requisitos, §5.2); Pohl (2010), escenarios y casos de uso; Cohn (2004), user stories y criterios de aceptación; SWEBOK v4.0 (KA Software Requirements); IREB CPRE FL v3.1 (UD~4, trabajo con requisitos en lenguaje natural y modelos).\\
\hline
\end{tabularx}

% ============================================================
% 2. OBJETIVOS
% ============================================================
\section*{\textcolor{utegreendk}{2. OBJETIVOS DE LA PRÁCTICA}}

\begin{tcolorbox}[colback=boxgray,colframe=utegreen,boxrule=0.8pt,arc=2mm]
\textbf{Objetivo general.} Modelar la funcionalidad del sistema del Proyecto Fin de Curso mediante un diagrama de casos de uso construido en una herramienta CASE profesional, y complementarlo con user stories documentadas con criterios de aceptación verificables, estableciendo la trazabilidad entre ambos artefactos y los requisitos elicitados.

\medskip
\textbf{Objetivos específicos.}
\begin{enumerate}[leftmargin=18pt,itemsep=2pt]
  \item Identificar los actores del sistema del PFC y delimitar la frontera del sistema (\textit{system boundary}).
  \item Construir un diagrama de casos de uso UML sintácticamente correcto, aplicando adecuadamente las relaciones \texttt{<<include>>}, \texttt{<<extend>>} y generalización.
  \item Especificar textualmente los casos de uso principales mediante una plantilla estructurada (flujo principal, flujos alternativos y excepciones).
  \item Redactar tres user stories conforme a la plantilla Connextra y a los criterios INVEST.
  \item Formular criterios de aceptación verificables para cada user story, en formato de lista de verificación o Dado--Cuando--Entonces (Gherkin).
  \item Trazar los casos de uso y las user stories hacia los requisitos elicitados del PFC, garantizando consistencia entre artefactos.
  \item Manejar una herramienta CASE (Enterprise Architect o StarUML) como entorno profesional de modelado y gestión de artefactos de requisitos.
\end{enumerate}
\end{tcolorbox}

% ============================================================
% 3. FUNDAMENTOS MÍNIMOS REQUERIDOS
% ============================================================
\section*{\textcolor{utegreendk}{3. FUNDAMENTOS MÍNIMOS REQUERIDOS}}

Antes de iniciar la práctica, el equipo debe dominar los siguientes conceptos, revisados en clase y disponibles en la bibliografía del curso:

\noindent\begin{tabularx}{\linewidth}{|L{0.30\linewidth}|X|}
\hline
\thdr{Concepto} & \thdr{Síntesis operativa}\\
\hline
\textbf{Actor} & Rol externo (persona, sistema u organización) que interactúa con el sistema. Los actores no forman parte del sistema: se ubican fuera de la frontera. Se distinguen actores primarios (inician el caso de uso) y secundarios (participan o reciben resultados).\\
\hline
\rowcolor{rowgray}\textbf{Caso de uso} & Secuencia de interacciones entre uno o más actores y el sistema que produce un resultado observable de valor para un actor. Se nombra con \textit{verbo en infinitivo + objeto} (p.~ej., \textit{Registrar pedido}).\\
\hline
\textbf{Relaciones UML} & \texttt{<<include>>}: comportamiento obligatorio reutilizado por el caso base. \texttt{<<extend>>}: comportamiento opcional o condicional que amplía el caso base en un punto de extensión. \textbf{Generalización}: especialización de un actor o caso de uso más general.\\
\hline
\rowcolor{rowgray}\textbf{Especificación textual} & Plantilla estructurada que documenta cada caso de uso: identificador, nombre, actores, precondiciones, flujo principal numerado, flujos alternativos, excepciones y postcondiciones.\\
\hline
\textbf{User story} & Escenario ágil breve en formato Connextra: \textit{Como} \textless rol\textgreater, \textit{quiero} \textless funcionalidad\textgreater\ \textit{para} \textless beneficio\textgreater. Debe cumplir los criterios INVEST: Independiente, Negociable, Valiosa, Estimable, Pequeña (\textit{Small}) y Testeable.\\
\hline
\rowcolor{rowgray}\textbf{Criterio de aceptación} & Condición verificable que determina cuándo la user story se considera satisfecha. Formatos admitidos: lista de verificación con condiciones medibles, o escenarios Gherkin (\textit{Dado--Cuando--Entonces}). Debe ser objetivo: no admite términos vagos sin cuantificar.\\
\hline
\end{tabularx}

% ============================================================
% 4. HERRAMIENTAS Y MATERIALES
% ============================================================
\section*{\textcolor{utegreendk}{4. HERRAMIENTAS Y MATERIALES}}

\begin{itemize}[leftmargin=18pt,itemsep=2pt]
  \item \textbf{Herramienta CASE (elegir una):}
    \begin{itemize}[leftmargin=14pt,itemsep=1pt]
      \item \textbf{Enterprise Architect} (Sparx Systems): versión de prueba o licencia académica disponible. Archivo de proyecto: \texttt{.qea}, \texttt{.eap} o \texttt{.eapx}.
      \item \textbf{StarUML} (MKLabs): versión de evaluación funcional. Archivo de proyecto: \texttt{.mdj}.
    \end{itemize}
  \item Documentos del PFC ya producidos: Entrega 1A (contexto, visión, alcance) e informe de elicitación (lista consolidada de requisitos con identificadores).
  \item Plantilla de especificación textual de casos de uso (sección 5, Fase 4 de esta guía).
  \item Repositorio GitHub del equipo, con la estructura de carpetas convenida para el PFC.
  \item Procesador de texto para el informe (exportación final en PDF).
\end{itemize}

\begin{tcolorbox}[colback=lightgold,colframe=gold,boxrule=0.8pt,arc=2mm,
  title=\textbf{Nota sobre herramientas},fonttitle=\bfseries]
No se admiten diagramas elaborados en herramientas de dibujo genérico (draw.io, Lucidchart, Canva, PowerPoint) ni imágenes generadas por terceros: el propósito de la práctica es el dominio de una herramienta CASE con repositorio de modelo. El archivo fuente nativo es evidencia obligatoria de autoría (ver criterios de admisión).
\end{tcolorbox}

% ============================================================
% 5. DESARROLLO DE LA PRÁCTICA
% ============================================================
\section*{\textcolor{utegreendk}{5. DESARROLLO DE LA PRÁCTICA (PROCEDIMIENTO DETALLADO)}}

\subsection*{\textcolor{midblue}{Fase 1 --- Preparación del entorno (estimado: 1 hora)}}
\begin{enumerate}[leftmargin=18pt,itemsep=2pt]
  \item Instalar la herramienta CASE seleccionada y verificar su correcto funcionamiento.
  \item Crear un proyecto nuevo con el nombre del sistema del PFC (p.~ej., \texttt{PFC\_NombreSistema\_EquipoX}).
  \item Crear dentro del proyecto un paquete (\textit{package}) denominado \texttt{Modelo de Casos de Uso} y, dentro de él, un diagrama de casos de uso vacío.
  \item Registrar capturas de pantalla del entorno configurado (estas capturas forman parte de las evidencias de autoría).
\end{enumerate}

\subsection*{\textcolor{midblue}{Fase 2 --- Identificación de actores y frontera del sistema (estimado: 1--2 horas)}}
\begin{enumerate}[leftmargin=18pt,itemsep=2pt]
  \item Revisar la lista consolidada de requisitos y la tabla de stakeholders del PFC.
  \item Derivar los \textbf{actores} del sistema: para cada stakeholder que interactúa directamente con el sistema, definir el rol que desempeña. Incluir, cuando corresponda, sistemas externos (pasarelas de pago, servicios de notificación, sistemas heredados).
  \item Elaborar una \textbf{tabla de actores} con tres columnas: nombre del actor, tipo (primario / secundario / sistema externo) y descripción de su interés u objetivo frente al sistema.
  \item Definir la \textbf{frontera del sistema}: qué funcionalidad queda dentro y qué queda fuera, en coherencia con el alcance declarado en la Entrega 1A.
\end{enumerate}

\subsection*{\textcolor{midblue}{Fase 3 --- Identificación y diagramación de casos de uso (estimado: 3--4 horas)}}
\begin{enumerate}[leftmargin=18pt,itemsep=2pt]
  \item A partir de los requisitos funcionales elicitados, identificar los casos de uso del sistema. Regla práctica: cada caso de uso debe producir un \textbf{resultado observable de valor} para un actor; las operaciones internas (p.~ej., \textit{validar contraseña}) no son casos de uso autónomos, sino pasos o inclusiones.
  \item Nombrar cada caso de uso con \textbf{verbo en infinitivo + objeto}, evitando nombres vagos (\textit{Gestionar sistema}) o tecnicismos de implementación (\textit{Insertar registro en BD}).
  \item Construir el diagrama en la herramienta CASE:
    \begin{itemize}[leftmargin=14pt,itemsep=1pt]
      \item Dibujar la frontera del sistema (rectángulo con el nombre del sistema).
      \item Ubicar los actores fuera de la frontera y los casos de uso dentro.
      \item Conectar actores y casos de uso mediante asociaciones.
      \item Aplicar \texttt{<<include>>} solo para comportamiento obligatorio compartido; \texttt{<<extend>>} solo para comportamiento condicional u opcional (con la dirección de la flecha correcta en cada relación); generalización solo cuando exista una verdadera relación de especialización.
    \end{itemize}
  \item Alcance mínimo del diagrama: \textbf{al menos 8 casos de uso}, \textbf{al menos 3 actores} y \textbf{al menos una relación de cada tipo} (\texttt{include}, \texttt{extend}) debidamente justificada. La generalización se incluye solo si el dominio la amerita; no debe forzarse.
  \item Asignar a cada caso de uso un \textbf{identificador único} (convención sugerida: \texttt{CU-01}, \texttt{CU-02}, \ldots) registrado como propiedad del elemento en la herramienta.
\end{enumerate}

\subsection*{\textcolor{midblue}{Fase 4 --- Especificación textual de casos de uso (estimado: 2--3 horas)}}
\begin{enumerate}[leftmargin=18pt,itemsep=2pt]
  \item Seleccionar los \textbf{tres casos de uso de mayor relevancia} para el sistema (los que concentran el valor central del PFC).
  \item Especificar cada uno con la siguiente plantilla:
\end{enumerate}

\noindent\begin{tabularx}{\linewidth}{|L{0.28\linewidth}|X|}
\hline
\bhdr{Campo} & \bhdr{Contenido}\\
\hline
Identificador & \texttt{CU-XX} (coincidente con el diagrama)\\
\hline
\rowcolor{rowgray}Nombre & Verbo en infinitivo + objeto\\
\hline
Actores & Primarios y secundarios involucrados\\
\hline
\rowcolor{rowgray}Precondiciones & Estado del sistema requerido antes de iniciar\\
\hline
Flujo principal & Secuencia numerada de pasos actor--sistema (escenario de éxito)\\
\hline
\rowcolor{rowgray}Flujos alternativos & Variaciones válidas del flujo principal, referenciadas al paso donde se bifurcan\\
\hline
Excepciones & Condiciones de error y respuesta del sistema\\
\hline
\rowcolor{rowgray}Postcondiciones & Estado del sistema tras la ejecución exitosa\\
\hline
Requisitos asociados & Identificadores de los requisitos del PFC que el caso de uso realiza\\
\hline
\end{tabularx}

\begin{enumerate}[leftmargin=18pt,itemsep=2pt,resume]
  \item Redactar los pasos del flujo principal alternando actor y sistema (\textit{1. El actor solicita\ldots\ 2. El sistema valida\ldots}), sin decisiones de diseño ni de interfaz prematuras.
  \item Verificar que cada especificación cumple las características de calidad de ISO/IEC/IEEE 29148: no ambigua, completa, singular, verificable y factible.
\end{enumerate}

\subsection*{\textcolor{midblue}{Fase 5 --- Documentación de tres user stories con criterios de aceptación (estimado: 2 horas)}}
\begin{enumerate}[leftmargin=18pt,itemsep=2pt]
  \item Seleccionar tres funcionalidades del sistema del PFC, asociadas a casos de uso del diagrama (puede tratarse de los mismos tres casos especificados en la Fase 4 o de otros, siempre que la trazabilidad quede explícita).
  \item Redactar cada user story con la plantilla Connextra completa:
  \begin{tcolorbox}[colback=lightblue,colframe=midblue,boxrule=0.6pt,arc=2mm]
    \textbf{US-XX.} \textit{Como} \textless rol concreto del sistema\textgreater, \textit{quiero} \textless funcionalidad específica\textgreater\ \textit{para} \textless beneficio o valor verificable\textgreater.
  \end{tcolorbox}
  \item Validar cada user story contra los criterios \textbf{INVEST}, dejando constancia escrita de la verificación (tabla o lista de chequeo por historia).
  \item Formular para cada user story \textbf{al menos tres criterios de aceptación} en uno de los dos formatos:
    \begin{itemize}[leftmargin=14pt,itemsep=1pt]
      \item \textbf{Gherkin}: \textit{Dado} \textless contexto\textgreater, \textit{cuando} \textless acción\textgreater, \textit{entonces} \textless resultado observable\textgreater.
      \item \textbf{Lista de verificación}: condiciones medibles y objetivas (incluyendo, donde aplique, valores límite, mensajes esperados y casos negativos).
    \end{itemize}
  \item Los criterios deben cubrir el \textbf{escenario de éxito}, al menos un \textbf{escenario alternativo o de borde} y al menos un \textbf{escenario de error}. Términos vagos sin métrica («rápido», «fácil», «seguro») invalidan el criterio.
  \item Registrar las user stories también dentro de la herramienta CASE (como elementos de requisito o notas vinculadas a los casos de uso), evidenciando la gestión de artefactos en el repositorio del modelo.
\end{enumerate}

\subsection*{\textcolor{midblue}{Fase 6 --- Trazabilidad y consistencia (estimado: 1 hora)}}
\begin{enumerate}[leftmargin=18pt,itemsep=2pt]
  \item Construir la \textbf{matriz de trazabilidad} con las columnas: requisito del PFC (\texttt{RF-XX}) $\rightarrow$ caso de uso (\texttt{CU-XX}) $\rightarrow$ user story (\texttt{US-XX}, cuando exista).
  \item Verificar consistencia: ningún caso de uso sin requisito de origen; ningún requisito funcional central sin caso de uso que lo realice (o con descarte justificado); las tres user stories trazadas a casos de uso del diagrama.
  \item Revisar la coherencia terminológica entre artefactos (mismos nombres de actores, roles y conceptos del dominio).
\end{enumerate}

\subsection*{\textcolor{midblue}{Fase 7 --- Elaboración del informe y entrega (estimado: 1--2 horas)}}
\begin{enumerate}[leftmargin=18pt,itemsep=2pt]
  \item Estructurar el informe en PDF con las siguientes secciones:
    \begin{enumerate}[leftmargin=16pt,itemsep=1pt,label=\alph*)]
      \item Portada (asignatura, equipo, integrantes, paralelo, fecha, sistema del PFC).
      \item Introducción breve: propósito del modelo y alcance cubierto.
      \item Tabla de actores y definición de la frontera del sistema.
      \item Diagrama de casos de uso exportado desde la herramienta CASE (imagen de alta resolución, legible).
      \item Justificación de cada relación \texttt{include}/\texttt{extend}/generalización empleada.
      \item Especificaciones textuales de los tres casos de uso seleccionados.
      \item Las tres user stories con su verificación INVEST y sus criterios de aceptación.
      \item Matriz de trazabilidad requisito $\rightarrow$ caso de uso $\rightarrow$ user story.
      \item Evidencias de uso de la herramienta (capturas del entorno con el proyecto abierto, estructura del repositorio del modelo).
      \item Conclusiones del equipo y distribución del trabajo entre integrantes.
    \end{enumerate}
  \item Subir al repositorio GitHub del equipo, en la carpeta convenida para esta práctica: (i) el informe PDF y (ii) el \textbf{archivo fuente nativo} del proyecto CASE. Verificar que el \textit{commit} se realice dentro del plazo.
\end{enumerate}

% ============================================================
% 6. CRITERIOS DE ADMISIÓN (GATEKEEPER)
% ============================================================
\section*{\textcolor{utegreendk}{6. CRITERIOS DE ADMISIÓN (GATEKEEPER)}}

\begin{tcolorbox}[colback=lightred,colframe=red!60!black,
  title=\textbf{CRITERIOS DE ADMISIÓN --- verificación previa a la revisión},
  fonttitle=\bfseries,boxrule=0.8pt,arc=2mm]
El incumplimiento de \textbf{cualquiera} de los siguientes criterios suspende la revisión detallada del trabajo y asigna la \textbf{nota mínima (25\,\% del puntaje, equivalente a nivel Insuficiente en todos los criterios)}. La decisión final de aplicación de la penalización corresponde al docente.
\begin{enumerate}[leftmargin=18pt,itemsep=2pt]
  \item \textbf{Archivo fuente nativo:} el repositorio contiene el archivo de proyecto de la herramienta CASE (\texttt{.qea}/\texttt{.eap}/\texttt{.eapx} o \texttt{.mdj}) y este abre correctamente, conteniendo el modelo presentado en el informe. No se admiten únicamente imágenes del diagrama.
  \item \textbf{Correspondencia con el PFC:} el modelo corresponde al sistema del Proyecto Fin de Curso del equipo (no a un caso genérico, de tutorial o de otro equipo) y es coherente con el alcance de la Entrega 1A y con los requisitos elicitados.
  \item \textbf{Evidencias de autoría:} el informe incluye capturas del entorno de la herramienta con el proyecto del equipo abierto, conforme a las disposiciones de autenticidad y antifraude del PFC.
  \item \textbf{Entrega en plazo y formato:} informe en PDF y archivo fuente subidos al repositorio GitHub del equipo en la carpeta y con la convención de nombres establecidas, dentro del plazo definido.
\end{enumerate}
\end{tcolorbox}

% ============================================================
% 7. ESCALA Y FÓRMULA DE CALIFICACIÓN
% ============================================================
\section*{\textcolor{utegreendk}{7. ESCALA Y FÓRMULA DE CALIFICACIÓN}}

\noindent\begin{tabularx}{\linewidth}{|C{0.22\linewidth}|C{0.16\linewidth}|X|}
\hline
\thdr{Nivel} & \thdr{Valor} & \thdr{Descripción general}\\
\hline
\cellcolor{lvlexc}\textbf{Excelente} & 4 (100\,\%) & Cumplimiento pleno del criterio, con calidad profesional y sin errores relevantes.\\
\hline
\cellcolor{lvlsat}\textbf{Satisfactorio} & 3 (75\,\%) & Cumplimiento sustancial con errores menores que no comprometen la validez del artefacto.\\
\hline
\cellcolor{lvldev}\textbf{En desarrollo} & 2 (50\,\%) & Cumplimiento parcial con errores conceptuales o de completitud que requieren corrección.\\
\hline
\cellcolor{lvlins}\textbf{Insuficiente} & 1 (25\,\%) & Cumplimiento mínimo, ausente o con errores graves que invalidan el artefacto.\\
\hline
\end{tabularx}

\begin{tcolorbox}[colback=boxgray,colframe=utegreen,boxrule=0.8pt,arc=2mm]
\textbf{Fórmula de calificación:}
\[
\text{Nota} \;=\; \left( \frac{\sum_{i=1}^{6} \text{Nivel}_i \times \text{Peso}_i}{4} \right) \times \text{Escala}
\]
donde \(\text{Nivel}_i \in \{1,2,3,4\}\) es el nivel alcanzado en el criterio \(i\), \(\text{Peso}_i\) su ponderación (la suma de pesos es 1) y \(\text{Escala}\) el puntaje máximo de la actividad definido en el plan de evaluación del curso (p.~ej., 10 puntos).
\end{tcolorbox}

% ============================================================
% 8. RÚBRICA ANALÍTICA (APAISADA)
% ============================================================
\begin{landscape}

\section*{\textcolor{utegreendk}{8. RÚBRICA ANALÍTICA DE EVALUACIÓN}}

{\small
\begin{longtable}{|L{3.4cm}|L{4.7cm}|L{4.7cm}|L{4.7cm}|L{4.7cm}|}
\hline
\thdr{Criterio (peso)} &
\cellcolor{lvlexc}\textbf{Excelente (4)} &
\cellcolor{lvlsat}\textbf{Satisfactorio (3)} &
\cellcolor{lvldev}\textbf{En desarrollo (2)} &
\cellcolor{lvlins}\textbf{Insuficiente (1)}\\
\hline
\endfirsthead
\hline
\thdr{Criterio (peso)} &
\cellcolor{lvlexc}\textbf{Excelente (4)} &
\cellcolor{lvlsat}\textbf{Satisfactorio (3)} &
\cellcolor{lvldev}\textbf{En desarrollo (2)} &
\cellcolor{lvlins}\textbf{Insuficiente (1)}\\
\hline
\endhead
\hline
\multicolumn{5}{r}{\footnotesize\textit{Continúa en la página siguiente\ldots}}\\
\endfoot
\hline
\endlastfoot

\textbf{C1. Actores y frontera del sistema}\newline\cellcolor{lightgold}(10\,\%) &
Tabla de actores completa (primarios, secundarios y sistemas externos cuando existen), con descripción del interés de cada uno; frontera del sistema explícita y coherente con el alcance de la Entrega 1A; ningún actor interno al sistema ni stakeholder sin interacción directa modelado como actor. &
Actores principales correctamente identificados con descripciones breves; frontera definida con desajustes menores frente al alcance del PFC; a lo sumo un actor cuestionable (rol redundante o secundario mal clasificado). &
Lista de actores incompleta (falta algún actor evidente del dominio) o con errores de concepto (componentes internos o stakeholders sin interacción modelados como actores); frontera implícita o inconsistente. &
Actores genéricos no derivados del PFC, ausentes o conceptualmente erróneos en su mayoría; sin frontera del sistema; tabla de actores ausente.\\
\hline

\rowcolor{rowgray}
\textbf{C2. Diagrama de casos de uso y relaciones UML}\newline\cellcolor{lightgold}(25\,\%) &
Diagrama con 8 o más casos de uso correctamente nombrados (verbo en infinitivo + objeto), identificadores únicos, asociaciones actor--caso correctas; relaciones \texttt{include} y \texttt{extend} aplicadas con semántica y dirección correctas y justificadas por escrito; sin casos de uso «funcionales internos» ni descomposición funcional tipo DFD; notación UML impecable. &
Diagrama con el alcance mínimo cumplido; nombres mayormente correctos; a lo sumo un error de semántica o dirección en \texttt{include}/\texttt{extend}, o una justificación débil; notación con detalles menores (estereotipos mal escritos, alguna asociación con flecha innecesaria). &
Diagrama por debajo del alcance mínimo (6--7 casos de uso) o con errores repetidos: relaciones \texttt{include}/\texttt{extend} confundidas entre sí o usadas como secuencia de pasos; nombres vagos o de implementación; varios casos de uso sin asociación a actor. &
Diagrama ausente, ilegible o con menos de 6 casos de uso; relaciones UML usadas sin semántica reconocible; descomposición funcional en lugar de casos de uso; diagrama no elaborado en la herramienta CASE.\\
\hline

\textbf{C3. Especificación textual de casos de uso}\newline\cellcolor{lightgold}(20\,\%) &
Tres especificaciones completas con todos los campos de la plantilla; flujos principales numerados alternando actor--sistema; flujos alternativos y excepciones referenciados al paso de bifurcación; pre/postcondiciones verificables; requisitos asociados identificados; redacción conforme a las características de calidad de ISO/IEC/IEEE 29148. &
Tres especificaciones con la plantilla completa y flujos principales correctos; flujos alternativos o excepciones presentes pero incompletos en una de las especificaciones; ambigüedades puntuales de redacción que no impiden la comprensión. &
Solo dos especificaciones completas, o tres con campos relevantes vacíos (excepciones, postcondiciones); flujos que mezclan decisiones de diseño o de interfaz; pasos sin alternancia actor--sistema; trazabilidad a requisitos omitida. &
Una o ninguna especificación utilizable; plantilla ignorada o reducida a una descripción narrativa libre; flujos inexistentes o incoherentes con el diagrama.\\
\hline

\rowcolor{rowgray}
\textbf{C4. User stories (formato y calidad INVEST)}\newline\cellcolor{lightgold}(20\,\%) &
Tres user stories en plantilla Connextra completa con rol concreto, funcionalidad específica y beneficio verificable; verificación INVEST documentada por historia con juicio argumentado; historias independientes entre sí, de tamaño implementable y claramente derivadas del PFC. &
Tres user stories en formato correcto; verificación INVEST presente aunque superficial en alguna historia; un beneficio genérico («para mejorar la experiencia») o un rol amplio en una de las historias. &
Tres historias con errores de formato (cláusula \textit{para} ausente, rol genérico tipo «usuario» en todas) o solo dos historias completas; verificación INVEST ausente o meramente declarativa (sin análisis); historias que son tareas técnicas y no escenarios de usuario. &
Menos de dos user stories utilizables; formato Connextra ignorado; historias copiadas de plantillas ajenas al PFC o indistinguibles de requisitos tradicionales sin rol ni beneficio.\\
\hline

\textbf{C5. Criterios de aceptación}\newline\cellcolor{lightgold}(15\,\%) &
Al menos tres criterios por historia, en Gherkin o lista de verificación, todos objetivos y medibles; cobertura del escenario de éxito, de al menos un caso alternativo o de borde y de al menos un caso de error por historia; valores límite o métricas explícitas donde corresponde; consistencia plena con la user story. &
Tres criterios por historia mayormente verificables; cobertura de éxito y error presente, con casos de borde débiles o ausentes en una historia; a lo sumo un criterio con redacción mejorable pero comprobable. &
Menos de tres criterios en alguna historia, o criterios redactados como descripción de la funcionalidad y no como condiciones verificables; términos vagos sin cuantificar («rápido», «intuitivo») en varios criterios; sin escenarios de error. &
Criterios de aceptación ausentes en una o más historias, o todos subjetivos e inverificables; criterios contradictorios con la user story; copia evidente de ejemplos genéricos.\\
\hline

\rowcolor{rowgray}
\textbf{C6. Trazabilidad, consistencia y manejo de la herramienta CASE}\newline\cellcolor{lightgold}(10\,\%) &
Matriz requisito~$\rightarrow$~caso de uso~$\rightarrow$~user story completa y sin elementos huérfanos (o con descartes justificados); terminología consistente entre artefactos; user stories registradas dentro del modelo CASE; identificadores gestionados como propiedades de los elementos; repositorio del modelo organizado en paquetes; evidencias de uso del entorno claras. &
Matriz de trazabilidad que cubre la mayoría de los elementos; consistencia terminológica general con discrepancias menores; modelo organizado en la herramienta, aunque las user stories consten solo en el informe; evidencias de uso suficientes. &
Trazabilidad parcial (varios casos de uso sin requisito de origen o user stories sin vínculo al diagrama); inconsistencias de nombres entre artefactos; modelo desorganizado en la herramienta (todo en un único paquete sin identificadores). &
Sin matriz de trazabilidad; artefactos desconectados entre sí; sin evidencia de gestión del modelo en la herramienta más allá del dibujo del diagrama.\\
\hline

\end{longtable}
}

\end{landscape}

% ============================================================
% 9. INTERPRETACIÓN DE RESULTADOS
% ============================================================
\section*{\textcolor{utegreendk}{9. INTERPRETACIÓN DE RESULTADOS}}

\noindent\begin{tabularx}{\linewidth}{|C{0.25\linewidth}|X|}
\hline
\thdr{Banda (\% del puntaje)} & \thdr{Interpretación}\\
\hline
\cellcolor{lvlexc}90--100\,\% & Dominio sobresaliente del modelado de casos de uso y de la documentación ágil; los artefactos pueden integrarse directamente a la Entrega 1B del PFC.\\
\hline
\cellcolor{lvlsat}75--89\,\% & Dominio satisfactorio; los artefactos requieren ajustes menores antes de su integración al PFC.\\
\hline
\cellcolor{lvldev}50--74\,\% & Dominio parcial; existen errores conceptuales (semántica de relaciones UML, verificabilidad de criterios) que deben corregirse con retroalimentación del docente.\\
\hline
\cellcolor{lvlins}$<$50\,\% & Dominio insuficiente; el equipo debe rehacer los artefactos centrales y se recomienda tutoría antes de la siguiente entrega del PFC.\\
\hline
\end{tabularx}

% ============================================================
% 10. REFERENCIAS
% ============================================================
\section*{\textcolor{utegreendk}{10. REFERENCIAS}}

\begin{itemize}[leftmargin=18pt,itemsep=2pt]
  \item ISO/IEC/IEEE. (2018). \textit{ISO/IEC/IEEE 29148:2018 --- Systems and software engineering --- Life cycle processes --- Requirements engineering}. ISO/IEC/IEEE.
  \item Pohl, K. (2010). \textit{Requirements Engineering: Fundamentals, Principles, and Techniques}. Springer.
  \item Cohn, M. (2004). \textit{User Stories Applied: For Agile Software Development}. Addison-Wesley.
  \item Bourque, P., \& Fairley, R. E. (Eds.). (2024). \textit{Guide to the Software Engineering Body of Knowledge (SWEBOK Guide), Version 4.0}. IEEE Computer Society.
  \item IREB. (2024). \textit{CPRE Foundation Level Syllabus, v3.1}. International Requirements Engineering Board.
  \item Object Management Group. (2017). \textit{OMG Unified Modeling Language (OMG UML), Version 2.5.1}. OMG.
\end{itemize}

\vspace{10pt}
\noindent\textcolor{gold}{\rule{\linewidth}{1.5pt}}
\begin{center}
\footnotesize Documento de uso académico --- Ingeniería de Requisitos (ISR-401) --- UTEQ, 2026--2027 PPA.
\end{center}

\end{document}
